\documentclass[conference]{IEEEtran}
\usepackage{amsmath,amssymb,booktabs,graphicx,url}
\usepackage[hidelinks]{hyperref}
\newcommand{\pending}{\textit{Pending}}
\title{A Reproducible Grayscale U-Net Baseline for Solar Filament Segmentation}
\author{\IEEEauthorblockN{Author name to be confirmed}
\IEEEauthorblockA{Affiliation and contact to be supplied}}
\begin{document}
\maketitle
\begin{abstract}
Solar filaments exhibit thin, irregular structures that make pixel segmentation
and separation of individual objects challenging. This preliminary technical
report describes a reproducible computer-vision baseline for the Solar Filament
Segmentation Challenge 2026. A grayscale U-Net with an ImageNet-initialized
ResNet18 encoder is trained using binary cross-entropy and soft Dice loss.
Annotation sets are preserved by annotator, while observing-day grouping
separates training, tuning, and holdout data. An intensity-contrast baseline
provides a classical comparison. Probabilities are restored to the original
image dimensions before connected-component extraction and compressed COCO
mask export. The first 20-epoch run completed in approximately 48 minutes as
reported by the experimenter. Numerical segmentation metrics and the accepted
competition score remain to be inserted from saved experiment outputs.
This draft makes no claim of state-of-the-art performance or architectural novelty.
\end{abstract}
\begin{IEEEkeywords}
solar filaments, H-alpha, semantic segmentation, U-Net, reproducibility
\end{IEEEkeywords}

\section{Introduction}
Reliable identification of filaments in solar H-alpha observations supports
the analysis of solar magnetic structures and eruptive activity.
Filaments are difficult targets because they are often faint, elongated,
and bounded by fine barbs. A method can achieve substantial foreground overlap
while incorrectly merging neighboring filaments or fragmenting one filament
into several predictions. Pixel overlap and object-level detection quality
must therefore be distinguished.

The MAGFiLO dataset provides manually annotated GONG observations with
filament segmentations, bounding boxes, spines, and chirality labels
\cite{magfilo}. The present work studies the competition's segmentation task,
not the prediction of chirality or solar eruptions. Our contribution is an
explicitly documented baseline and evaluation/export procedure, rather than a
new neural architecture.

\section{Related Work and Research Context}
MAGFiLO v1.0 reports 10,244 annotated filaments from 1,593 observations and
inter-annotator agreement of $\kappa=0.66$ \cite{magfilo}. These published dataset
totals should not be confused with the physical JPEG counts in the competition
release or with annotator-specific image records.

The supplied ESAIR literature addresses two complementary problems.
H-Alpha Anomalyzer detects corrupt observations using regional mean-intensity
statistics and explainable grid-level anomaly scores \cite{softwarex,icdmw}.
Its image-level anomaly F1 is not a segmentation score. Supervised SVM and
one-class SVM in the comparative study are quality-control classifiers,
not the supervised filament segmenter implemented here.

The 2026 ESAIR poster describes U-Net segmentation with multiscale IoU,
geometry-based fragment grouping, and ConvNeXt classification of extracted
filament regions \cite{poster}. We do not reproduce the poster's loss or
grouping algorithm: their complete implementation details are not present
in the supplied poster. Our baseline instead uses an established U-Net
\cite{unet} and connected components. Bounding boxes are derived from masks;
no separate box detector or bounding-box regression model is trained.

\section{Data and Experimental Protocol}
The reviewed competition release contains 707 training JPEGs and 180
unlabeled test JPEGs. Training annotations contain 1,154 annotator-specific
image records and 8,199 filament annotations. Images are $2048\times2048$
8-bit grayscale observations. Multiple annotator records can refer to the
same photograph.

Each annotator-specific annotation set is retained as a separate supervised
target. Its filament masks are unioned into a binary foreground target.
Different annotators' masks are not merged into consensus labels. All
chirality categories are treated as foreground. Polygons are decoded using
COCO-compatible mask operations, rather than interpreting the JSON as a
filename-to-mask dictionary.

Observing days are assigned to disjoint training, tuning, and holdout
partitions with seed 2026. For the inspected release this yields 499, 103,
and 105 physical images, corresponding to 827, 161, and 166 annotation sets.
Same-day grouping prevents duplicate photographs and same-day observations
from crossing partitions; it does not guarantee independence of long-lived
events across adjacent days. Multiply annotated images receive multiple
training exposures, a limitation of the chosen target policy.

\section{Methods}
\subsection{Grayscale neural segmentation}
Inputs are resized bilinearly to $512\times512$. Binary targets use nearest
neighbor resizing. With grayscale value $v_i$, the input normalization is
$x_i=2(v_i/255)-1$. Paired quarter-turn rotations and horizontal flips are
applied during training.

The default network is U-Net with a ResNet18 encoder \cite{resnet}, not FPN.
ImageNet weights initialize the encoder, adapted by the segmentation library
to one input channel. Hidden convolutions and nonlinearities remain present.
The output head produces logits $z_i$; setting its activation to \texttt{None}
does not remove hidden activations. Probabilities are
\begin{equation}
p_i=\sigma(z_i)=(1+\exp(-z_i))^{-1}.
\end{equation}
This is supervised binary semantic segmentation, not continuous-target
regression.

For binary reference $y_i$, the loss is
\begin{equation}
L=\tfrac12 L_{\mathrm{BCE}}+\tfrac12 L_{\mathrm{Dice}},
\end{equation}
\begin{equation}
L_{\mathrm{BCE}}=-\frac{1}{N}\sum_i
[y_i\log p_i+(1-y_i)\log(1-p_i)],
\end{equation}
\begin{equation}
L_{\mathrm{Dice}}=1-\frac{1}{B}\sum_{b=1}^{B}
\frac{2\sum_i p_{bi}y_{bi}+1}{\sum_i p_{bi}+\sum_i y_{bi}+1}.
\end{equation}
The implementation evaluates BCE directly from logits for stability.
AdamW uses learning rate $3\times10^{-4}$ and batch size eight.
The baseline trains for 20 epochs, selecting the checkpoint with the lowest
tuning loss. It has no baseline early stopping or learning-rate scheduler.

\subsection{Classical comparator and instance extraction}
The contrast comparator estimates a smoothed background
$B=G_\sigma*I$ and selects dark departures within an eroded heuristic disk mask:
\begin{equation}
M(x)=D(x)\land[B(x)-I(x)>\tau_c].
\end{equation}
It is a tuned heuristic rather than a trained SVM or a reproduction of
Anomalyzer. Sunspots, limb artifacts, and clouds can confound it.

Neural probabilities are bilinearly restored to $2048\times2048$ before
thresholding. Eight-connected components define predicted instances.
Probability thresholds $\{0.35,0.50,0.65\}$ and minimum original-resolution
areas $\{64,256\}$ are selected using tuning PQ on at most 32 physical tuning
images. The contrast baseline searches thresholds $\{0.025,0.04,0.06\}$ with
the same area choices. These values are search candidates, not demonstrated
optimal settings.

\subsection{Evaluation and artifacts}
For foreground pixel counts,
\begin{equation}
\mathrm{IoU}=\frac{TP}{TP+FP+FN},\qquad
\mathrm{Dice}=\frac{2TP}{2TP+FP+FN}.
\end{equation}
Binary foreground Dice equals pixel F1 for the same masks and conventions.
Object-level PQ \cite{pq} uses one-to-one matches at IoU strictly greater
than 0.5:
\begin{equation}
\mathrm{PQ}=
\frac{\sum_{(p,g)\in\mathcal{M}}\mathrm{IoU}(p,g)}
{|\mathcal{M}|+\tfrac12 FP_{\mathrm{obj}}+\tfrac12 FN_{\mathrm{obj}}}.
\end{equation}
Local evaluation ignores chirality and averages annotator-specific scores
within each physical image before averaging images. Split/merge diagnostics
use overlap greater than 0.1. Local PQ is not asserted to reproduce every
private competition evaluator edge case.

The notebook saves weights, annotation and checkpoint hashes, split manifests,
calibration candidates, loss history, runtime, figures, and installed packages.
Video galleries display human outlines where available and predicted masks
and mask-derived boxes. Frames span different dates and stations; they do not
constitute a continuous solar evolution sequence.

\section{Preliminary Results and Planned Comparisons}
\begin{table}[t]
\centering
\caption{Baseline evidence status. Pending entries must be replaced with
actual saved outputs before submission. No leaderboard score is inferred.}
\begin{tabular}{ll}
\toprule
Quantity & Status/value\\
\midrule
Training budget & 20 epochs\\
Reported run duration & approximately 48 min\\
Selected checkpoint epoch & \pending\\
Holdout PQ / Dice / IoU & \pending\\
Contrast holdout PQ & \pending\\
Accepted leaderboard score & \pending\\
\bottomrule
\end{tabular}
\end{table}
The experimenter reported approximately 48 minutes for the first run; the
timed scope and hardware must be checked against runtime artifacts.
The initial test export reported zero retained instances for three of 180
images. This may indicate missed faint filaments, area filtering, or genuinely
empty scenes; absent test labels, the cause cannot be established from counts.
No artificial filament masks should be inserted to conceal this outcome.
Candidate export omits filament rows for zero-detection images and records
them explicitly. Evaluator acceptance of that representation remains to be
confirmed by the actual competition response.

Two optional follow-up experiments are implemented but not yet measured:
fresh ImageNet initialization for at most 30 epochs at $3\times10^{-4}$, and
fine-tuning the baseline best checkpoint for at most ten new epochs at
$10^{-4}$. Both use a new optimizer and early stopping after five consecutive
epochs without tuning-loss improvement greater than $10^{-4}$. Fine-tuning
retains its starting checkpoint as a candidate. It is not exact optimizer-state
resume from epoch 20. Both are calibrated on the same tuning subset; neither
should be selected using hidden-test feedback or repeated holdout inspection.

\section{Discussion and Limitations}
Correct annotation mapping and full-resolution export are necessary for a
meaningful baseline. The supplied reference notebooks use a direct filename
lookup on a COCO top-level dictionary; this returns empty targets for the
reviewed release. Their saved falling losses alone therefore do not establish
filament learning. Their saved inference paths also encode resized masks
without restoring the required canvas. This observation concerns the supplied
source versions, not independently verified leaderboard rankings.

Our implementation repairs those data-flow issues but does not establish
superior accuracy. Downsampling can erase thin structures, and upsampling
cannot recover information already lost. Binary supervision discards instance
identity during learning. Connected components can merge touching filaments
or split fragmented ones. A single split and seed do not establish statistical
significance. Future work should isolate resolution, boundary losses,
geometry grouping, quality-aware weighting, and architecture choices using
controlled grouped experiments.

\section{Conclusion}
We describe an established-model grayscale segmentation baseline with explicit
annotation handling, grouped evaluation, auditable calibration, and full-size
mask export. Quantitative conclusions await insertion of measured holdout
results and accepted competition scores. The final report must identify the
actual selected weights, environment, runtime scope, and failure cases rather
than attributing performance to architecture reputation alone.

\section*{Reproducibility and draft checklist}
Before submitting: insert confirmed authors; use the organizer's exact
template and verify its four-page limit; insert actual metrics, hardware and
figures; resolve zero-detection evaluator acceptance; add a public repository
and video link; include weights and exact dependencies. The Overleaf project
was restricted during draft preparation. This source is a portable IEEE-style
draft, not a confirmed rendering in that project.

\bibliographystyle{IEEEtran}
\bibliography{references}
\end{document}
